\section{Data}

\subsection{Marina Bay Sands Review Dataset}

We collected 10,232 English-language reviews of Marina Bay
Sands, Singapore from TripAdvisor, spanning January 2014 to
October 2022. We filtered for four- and five-star reviews, targeting
reviews likely to contain improvement suggestions embedded within
positive text. After removing null entries, 8,571
reviews remained: 2,414 four-star and 6,157 five-star.

We segmented reviews into sentences using the spaCy
\texttt{en\_core\_web\_sm} model \citep{honnibal2020spacy} and
applied three sequential cleaning steps to the 72,391 segmented
sentences: (1)~a language filter removing 1 non-English sentence,
(2)~exact deduplication removing 1,137 duplicates, and (3)~a
length filter removing 1,168 sentences outside the 4--100 token
range. After all preprocessing, 70,085
sentences remained, with a median length of 15 tokens and a mean
of 8.4 sentences per review. Table~\ref{tab:mbs-stats}
summarises the dataset statistics.

\begin{table}[t]
\centering\small
\caption{MBS dataset statistics.}
\label{tab:mbs-stats}
\begin{tabular}{lr}
\toprule
\textbf{Statistic} & \textbf{Value} \\
\midrule
Raw reviews            & 10,232 \\
After 4--5 star filter & 8,571  \\
Four-star reviews      & 2,414  \\
Five-star reviews      & 6,157  \\
Sentences after preprocessing & 70,085 \\
Median sentence length (tokens) & 15 \\
Mean sentences per review & 8.4 \\
\bottomrule
\end{tabular}
\end{table}

\subsection{SemEval-2019 Task 9 Dataset}

We use the publicly available SemEval-2019 Task~9 Subtask~A dataset
\citep{negi2019semeval}, collected from the Uservoice software feedback forum.
The SemEval training set contains 8,500 sentences (2,085
suggestions, 6,415 non-suggestions; 24.5\% suggestion rate). The
SemEval test set contains 833 sentences (87 suggestions, 746
non-suggestions; 10.4\% suggestion rate). This class imbalance,
particularly in the SemEval test set, motivates our use of
suggestion-class
F1 rather than accuracy as the primary evaluation metric.

The SemEval dataset differs from our MBS sentences in several respects. Software
forum posts tend to be explicit feature requests using modal verbs (``The system
should allow\ldots''), while hotel review suggestions are often implicit,
embedded within descriptive or comparative language (``I wish the breakfast
had\ldots''). Sentence lengths and vocabulary also differ: the SemEval data
contains technical terms absent from hospitality reviews. These differences
constitute the domain shift that our cross-domain evaluation measures.

\begin{table}[t]
\centering\small
\caption{SemEval-2019 Task~9 dataset statistics.}
\label{tab:semeval-stats}
\begin{tabular}{lrrrr}
\toprule
\textbf{Split} & \textbf{Total} & \textbf{Sugg.} & \textbf{Non-sugg.} & \textbf{Rate} \\
\midrule
Training & 8,500 & 2,085 & 6,415 & 24.5\% \\
Test     &   833 &    87 &   746 & 10.4\% \\
\bottomrule
\end{tabular}
\end{table}

\subsection{Annotation}

We developed annotation guidelines adapted from the SemEval-2019
task definition \citep{negi2019semeval,negi2018open}, with
extensions for hotel-domain edge cases such as comparative
suggestions, implicit improvement requests, and the distinction
between complaints and concrete suggestions.

\paragraph{Calibration round.} All four annotators independently
  labelled the same MBS calibration sample of 100 sentences. To
  increase the proportion of informative examples, we used signal
  enrichment: 50\% of the sample was drawn from sentences
  containing suggestion-signal keywords (modal verbs, imperatives,
  conditionals), and 50\% was drawn randomly. Inter-annotator
  agreement was fair: Fleiss' $\kappa$ = 0.379. Pairwise Cohen's
  $\kappa$ ranged from 0.212 to 0.672 (mean 0.426). The highest
  pairwise agreement (0.672) occurred between two annotators who
  most consistently distinguished complaints from suggestions; the
  lowest (0.212) reflected differing thresholds for implicit
  suggestions. The moderate overall agreement is consistent with
  the subjective nature of suggestion boundaries reported by
  \citet{negi2018open}.

\paragraph{Disagreement resolution.} Sentences with majority
  agreement (3 or 4 of 4 annotators) received the majority label.
  Of the 100 calibration sentences, 11 produced 2-vs-2 ties; these
  were resolved by group discussion. Most ties involved sentences
  where negative sentiment co-occurred with modal verbs (e.g.\
  ``I would definitely stay again''), which annotators initially
  split on. After calibration, guidelines were revised to clarify
  two recurring edge cases: complaints without actionable direction
  and informational statements versus guest-directed advice.

\paragraph{Split annotation.} After calibration, 300 additional
  sentences were divided into four batches of 75, with each
  annotator labelling one batch independently. Combined with the
  100 calibration sentences, the MBS annotated set contains 400
  sentences: 73 suggestions (18.25\%) and 327 non-suggestions
  (81.75\%). The suggestion rate is lower than the SemEval training
  set (24.5\%), reflecting the less directive language typical of
  hotel reviews.

\paragraph{Train/test split.} The MBS annotated set was split
  80/20 (stratified by label) into an MBS train split of 320
  sentences (58 suggestions, 18.1\%) and an MBS test split of 80
  sentences (15 suggestions, 18.8\%).
